% chapters/ch04.tex
\chapter{Delimitaciones, alcances e hipótesis/supuestos de investigación}

\begin{refsection}
\addbibresource{references/ch04.bib}

\section{Delimitaciones de la investigación}

Las delimitaciones permiten acotar el alcance del estudio, evitar interpretaciones indebidas de los resultados y proteger la validez científica de la investigación. En investigaciones educativas complejas, especialmente aquellas mediadas por tecnología, la definición clara de los límites del estudio resulta fundamental para comprender el contexto en el que se generan los resultados y las condiciones bajo las cuales estos pueden ser interpretados o transferidos.

En el presente estudio, las delimitaciones se establecen en los planos temático, metodológico, poblacional y tecnológico.

\subsection{Delimitación temática}

La investigación se delimita al estudio del diagnóstico educativo inclusivo mediado por tecnología, con especial énfasis en el diseño y desarrollo de un sistema tecnológico que permita apoyar la identificación de necesidades educativas diversas a partir del análisis de datos educativos.

En este sentido, el estudio:

\begin{itemize}
\item No se centra en la implementación completa de reformas curriculares.
\item No evalúa políticas educativas a nivel macro o sistémico.
\item No pretende sustituir la labor docente mediante automatización tecnológica.
\end{itemize}

El foco principal se orienta hacia:

\begin{itemize}
\item el diseño conceptual y metodológico de un sistema tecnológico de diagnóstico educativo inclusivo;
\item el análisis de su potencial como instrumento de recopilación y análisis de datos educativos;
\item y su contribución al apoyo de la toma de decisiones pedagógicas orientadas a la inclusión educativa.
\end{itemize}

Esta delimitación temática permite concentrar la investigación en la intersección entre inclusión educativa, tecnología educativa y educación basada en datos.

\subsection{Delimitación metodológica}

Desde el punto de vista metodológico, la investigación se delimita a la aplicación de un enfoque mixto, combinando métodos cuantitativos y cualitativos dentro de un diseño de Investigación Basada en el Diseño (\DBR).

Este enfoque permite:

\begin{itemize}
\item analizar datos educativos generados por el sistema tecnológico;
\item interpretar las experiencias y percepciones de los actores educativos;
\item y mejorar progresivamente el diseño del sistema mediante ciclos iterativos de implementación y evaluación.
\end{itemize}

No se pretende realizar:

\begin{itemize}
\item experimentación clínica controlada;
\item estudios longitudinales a gran escala;
\item validaciones psicométricas exhaustivas de todos los instrumentos educativos disponibles.
\end{itemize}

En consecuencia, la investigación prioriza el desarrollo, implementación y evaluación contextualizada de una innovación educativa mediada por tecnología.

\subsection{Delimitación poblacional}

La población de estudio se delimita a estudiantes y docentes pertenecientes a contextos educativos específicos seleccionados bajo criterios de accesibilidad, pertinencia y diversidad.

La selección de los participantes se realiza mediante muestreo intencional, priorizando la inclusión de diferentes perfiles educativos y contextos de aprendizaje.

En este sentido, los resultados obtenidos:

\begin{itemize}
\item no pretenden generalizarse estadísticamente a toda la población educativa;
\item pero sí poseen potencial de transferibilidad conceptual y metodológica a contextos educativos similares.
\end{itemize}

Este enfoque resulta coherente con el carácter contextualizado de las investigaciones educativas y con la lógica iterativa propia del enfoque \DBR.

\subsection{Delimitación tecnológica}

La investigación se centra en el diseño y desarrollo de un sistema tecnológico prototipo funcional orientado al diagnóstico educativo inclusivo.

Este sistema ha sido concebido como una herramienta académica y de investigación, cuyo objetivo principal es explorar las posibilidades de la analítica de datos educativos para apoyar la identificación de necesidades de aprendizaje diversas.

En este sentido, el estudio:

\begin{itemize}
\item no pretende desarrollar un producto comercial;
\item no evalúa la viabilidad económica o de mercado del sistema;
\item ni busca su implementación a gran escala en sistemas educativos nacionales.
\end{itemize}

La delimitación tecnológica permite enfocar el análisis en el potencial pedagógico y metodológico de los sistemas inteligentes aplicados a la educación inclusiva.

\section{Alcances de la investigación}

Los alcances de la investigación describen las contribuciones que el estudio pretende aportar en los planos teórico, metodológico, práctico y social.

\subsection{Alcance teórico}

En el plano teórico, la investigación contribuye al fortalecimiento del marco conceptual sobre educación inclusiva mediada por tecnología, particularmente mediante la articulación entre tres enfoques fundamentales:

\begin{itemize}
\item inclusión educativa,
\item Diseño Universal para el Aprendizaje (\DUA),
\item educación basada en datos.
\end{itemize}

Esta articulación permite ampliar la comprensión del diagnóstico educativo desde una perspectiva multidimensional que reconoce la diversidad como una característica inherente de los procesos de aprendizaje.

Asimismo, el estudio aporta a la reflexión sobre el papel de los sistemas tecnológicos inteligentes como herramientas de apoyo al conocimiento pedagógico y a la toma de decisiones educativas.

\subsection{Alcance metodológico}

Desde una perspectiva metodológica, la investigación propone un modelo replicable para el diagnóstico educativo inclusivo basado en datos.

Entre los aportes metodológicos del estudio se encuentran:

\begin{itemize}
\item la aplicación del enfoque \DBR\ en el diseño y evaluación de sistemas tecnológicos educativos;
\item la integración de métodos cuantitativos y cualitativos en el análisis de datos educativos;
\item el desarrollo de instrumentos digitales para la recopilación sistemática de información educativa.
\end{itemize}

Estos aportes contribuyen a fortalecer las metodologías de investigación aplicadas al estudio de innovaciones educativas mediadas por tecnología \parencite{creswell2014,wang2005}.

\subsection{Alcance práctico y tecnológico}

Desde una perspectiva aplicada, la investigación propone un sistema tecnológico que puede ser utilizado como herramienta de apoyo para docentes e instituciones educativas.

Entre sus posibles aplicaciones se encuentran:

\begin{itemize}
\item la identificación temprana de necesidades educativas diversas;
\item el análisis de patrones de aprendizaje a partir de datos educativos;
\item la generación de recomendaciones pedagógicas alineadas con los principios del \DUA.
\end{itemize}

Este enfoque permite fortalecer los procesos de toma de decisiones pedagógicas fundamentadas en evidencia.

\subsection{Alcance social y educativo}

En el plano social y educativo, el estudio contribuye a promover prácticas educativas más equitativas e inclusivas mediante el uso responsable de tecnologías basadas en datos.

En particular, la investigación aporta a:

\begin{itemize}
\item la reducción de barreras para el aprendizaje;
\item el reconocimiento de la diversidad como valor educativo;
\item el desarrollo de prácticas pedagógicas fundamentadas en principios de equidad y accesibilidad.
\end{itemize}

Estos aportes se alinean con los marcos internacionales que promueven la inclusión educativa y el uso ético de la tecnología en educación \parencite{unesco2020}.

\section{Hipótesis y supuestos de la investigación}

Dado el enfoque mixto adoptado en la investigación, se articulan hipótesis (relacionadas con el componente cuantitativo) y supuestos de investigación (vinculados con la interpretación cualitativa y conceptual del fenómeno educativo).

Esta articulación permite analizar tanto relaciones observables entre variables como comprender las dinámicas educativas desde una perspectiva contextual.

\subsection{Hipótesis de investigación (componente cuantitativo)}

\subsubsection{Hipótesis general}

\textbf{H1:} El uso de un sistema tecnológico de diagnóstico educativo inclusivo mejora la identificación de necesidades educativas diversas y apoya la toma de decisiones pedagógicas fundamentadas en comparación con métodos diagnósticos tradicionales.

\subsubsection{Hipótesis específicas}

\begin{itemize}
\item \textbf{H1.1:} El sistema tecnológico permite identificar patrones de aprendizaje relevantes mediante el análisis de datos educativos.
\item \textbf{H1.2:} Los diagnósticos generados por el sistema presentan coherencia con los principios del \DUA.
\item \textbf{H1.3:} Los docentes perciben el sistema como una herramienta útil para apoyar prácticas educativas inclusivas.
\end{itemize}

\subsection{Supuestos de investigación (componente cualitativo y conceptual)}

Los supuestos de investigación orientan la interpretación del fenómeno educativo desde una perspectiva comprensiva.

En el contexto de este estudio, se asumen los siguientes principios:

\begin{itemize}
\item El diagnóstico educativo inclusivo requiere integrar múltiples fuentes de información.
\item La tecnología puede actuar como mediadora del conocimiento pedagógico, pero no sustituye el juicio profesional del docente.
\item El uso ético de los datos educativos constituye una condición fundamental para el desarrollo de sistemas tecnológicos inclusivos.
\item La diversidad en el aprendizaje no representa una desviación del proceso educativo, sino una característica inherente de los sistemas educativos.
\end{itemize}

\section{Articulación entre delimitaciones, alcances e hipótesis}

La coherencia metodológica de una investigación científica depende de la relación consistente entre los límites del estudio, los aportes esperados y las proposiciones que orientan el análisis empírico.

En este sentido, las delimitaciones, los alcances y las hipótesis constituyen componentes interrelacionados que permiten estructurar de manera rigurosa el diseño investigativo.

Las delimitaciones establecen el marco dentro del cual se desarrolla el estudio, definiendo su alcance temático, metodológico, poblacional y tecnológico. Los alcances describen las contribuciones potenciales del estudio en los planos teórico, metodológico, práctico y social. Finalmente, las hipótesis y los supuestos orientan el análisis empírico y permiten evaluar el grado en que la investigación contribuye al conocimiento científico.

La relación entre estos componentes del diseño investigativo se representa en la Figura \ref{fig:coherencia-investigacion}.

\begin{figure}[htbp]
\centering
\resizebox{\textwidth}{!}{%
\begin{tikzpicture}[
font=\footnotesize,
node distance=14mm,
block/.style={draw, rounded corners=2mm, align=center, inner sep=8pt, text width=4.5cm},
bigblock/.style={draw, rounded corners=2mm, align=center, inner sep=10pt, text width=8cm},
arrow/.style={-{Stealth[length=2.6mm]}, thick}
]

\node[block] (delim)
{\textbf{Delimitaciones}\\
Temáticas\\
Metodológicas\\
Poblacionales\\
Tecnológicas};

\node[block, right=30mm of delim] (alcances)
{\textbf{Alcances}\\
Teóricos\\
Metodológicos\\
Prácticos\\
Sociales};

\node[block, below=18mm of delim] (hipotesis)
{\textbf{Hipótesis y supuestos}\\
Hipótesis cuantitativas\\
Supuestos interpretativos};

\node[bigblock, below=18mm of alcances] (coherencia)
{\textbf{Coherencia del diseño investigativo}\\
Rigor metodológico\\
Claridad conceptual\\
Solidez científica};

\draw[arrow] (delim) -- (alcances);
\draw[arrow] (delim) -- (hipotesis);
\draw[arrow] (alcances) -- (coherencia);
\draw[arrow] (hipotesis) -- (coherencia);

\end{tikzpicture}
}

\caption{Relación entre delimitaciones, alcances e hipótesis dentro del diseño metodológico de la investigación.}
\label{fig:coherencia-investigacion}
\end{figure}

Como se observa en la Figura \ref{fig:coherencia-investigacion}, las delimitaciones establecen el marco dentro del cual se desarrolla la investigación, mientras que los alcances definen las contribuciones esperadas del estudio. Las hipótesis y los supuestos orientan el análisis empírico y permiten evaluar el grado en que la propuesta investigativa contribuye al conocimiento científico y a la mejora de las prácticas educativas inclusivas.

\printbibliography[heading=subbibliography,title={Referencias (Capítulo IV)}]

\end{refsection}